% ============================================================
% results.tex
% ============================================================

This chapter presents the full empirical results of our
competing-risks survival analysis. We proceed from
nonparametric characterization of litigation duration and
outcomes, through progressively richer covariate-adjusted
models, to robustness checks and machine learning
benchmarking. The organizing question throughout is: how do
the PSLRA, circuit geography, and case-level structural
characteristics jointly shape the timing and probability of
settlement versus dismissal?

\section{Nonparametric Duration and Outcome Estimates}

\subsection{Overall Duration}

Figure~\ref{fig:km_overall} presents the Kaplan-Meier
estimate of the overall survival function for the full
analysis sample of 12,968 federal securities class actions.
The curve drops steeply in the first two years, reflecting
the concentration of dismissals during the
motion-to-dismiss phase: approximately 25\% of cases are
resolved within the first year, and roughly 50\% within two
years. Resolution slows markedly after year four, with fewer
than 10\% of cases remaining pending beyond eight years.

\begin{figure}[htb]
    \centering
    \begin{subfigure}[b]{0.49\textwidth}
        \centering
        \includegraphics[width=\textwidth]{./Figures/figure1_km_overall.png}
        \caption{Kaplan-Meier Survival Function}
        \label{fig:km_overall}
    \end{subfigure}
    \hfill
    \begin{subfigure}[b]{0.49\textwidth}
        \centering
        \includegraphics[width=\textwidth]{./Figures/figure2_cif_overall.png}
        \caption{Aalen-Johansen CIF}
        \label{fig:cif_overall}
    \end{subfigure}
    \vskip 10pt
    \caption[Overall Duration and Competing-Outcome Patterns]{Overall duration and competing-outcome patterns. Left: Kaplan-Meier overall survival ($n = 12,968$); shaded band = 95\% pointwise CI. Right: Aalen-Johansen CIFs for settlement and dismissal under Scheme A.}
    \label{fig:baseline_plots}
\end{figure}

\subsection{Competing-Risks Cumulative Incidence: Overall}

Figure~\ref{fig:cif_overall} decomposes overall resolution
into its two competing terminal events. Dismissal incidence
climbs rapidly: by two years post-filing, roughly 35\% of
cases have been dismissed, and by five years the cumulative
dismissal probability approaches 60\%. Settlement incidence
accumulates slowly and reaches only approximately 20\% by
eight years. The two curves do not converge, reflecting the
fundamental asymmetry between the two outcomes: dismissal
is the dominant resolution pathway, operating primarily
through early pleading-stage motions, while settlement is
a comparatively rare and temporally dispersed event
concentrated among cases that survive the initial
motion-to-dismiss.

A naive Kaplan-Meier estimate of time to settlement that
censors dismissals would produce upwardly biased settlement
probabilities by treating dismissed cases as though they
remained at risk of eventual settlement---precisely the
type of bias the Aalen-Johansen competing-risks estimator
avoids.

\subsection{Cumulative Incidence by PSLRA Regime}

Figure~\ref{fig:cif_pslra} stratifies the CIFs by PSLRA
regime. Post-PSLRA dismissal incidence rises sharply in
the first year and reaches nearly 63\% by five years,
substantially above the pre-PSLRA trajectory of 42\%.
Post-PSLRA settlement incidence is simultaneously
suppressed, reaching only 15.7\% at five years compared
with 43.8\% pre-PSLRA. The reform affects both competing
processes simultaneously and in opposite directions.

\begin{figure}[htb]
    \centering
    \includegraphics[width=\linewidth]{./Figures/figure3_cif_pslra.png}
    \vskip 5pt
    \caption[Cumulative Incidence of Outcomes by PSLRA Regime]{Aalen-Johansen CIFs stratified by PSLRA regime. Solid lines = Pre-PSLRA ($n = 1,032$); dashed = Post-PSLRA ($n = 11,936$). Green = settlement; orange/red = dismissal. Gray's test rejects equality for both settlement ($\text{stat} = 397.6, p < 0.001$) and dismissal ($\text{stat} = 179.8, p < 0.001$).}
    \label{fig:cif_pslra}
\end{figure}

Table~\ref{tab:cif_horizons} quantifies differences at four
key horizons. The PSLRA effect is particularly pronounced
in the first year: one-year dismissal incidence is 29.6\%
post-PSLRA versus 15.0\% pre-PSLRA. By contrast, the
settlement gap narrows only modestly over time: the five-year
settlement difference is 28.1 percentage points, indicating
that the PSLRA's suppression of settlement is a permanent
reduction in outcome probability rather than a timing
effect. Notably, pre-PSLRA five-year dismissal and
settlement CIFs are nearly equal (42.2\% vs.\ 43.8\%),
suggesting that under the old regime the two outcomes were
roughly equiprobable for long-duration cases. Post-PSLRA,
dismissal probability far exceeds settlement probability
at every horizon.

\begin{table}[htbp]
\centering
\caption{Cumulative Incidence at Key Horizons by PSLRA Regime (\%), Aalen-Johansen Estimates}
\label{tab:cif_horizons}
\renewcommand{\arraystretch}{1.3} % Adds vertical padding for a professional feel
\begin{tabular*}{\textwidth}{@{\extracolsep{\fill}} l l r r r r @{}}
\toprule
\textbf{Regime} & \textbf{Outcome} & \textbf{1 yr} & \textbf{2 yr} & \textbf{3 yr} & \textbf{5 yr} \\
\midrule
Pre-PSLRA  & Settlement &  4.7 & 19.1 & 28.1 & 43.8 \\
           & Dismissal  & 15.0 & 25.9 & 29.5 & 42.2 \\
\midrule
Post-PSLRA & Settlement &  2.0 &  5.8 & 10.1 & 15.7 \\
           & Dismissal  & 29.6 & 38.9 & 48.5 & 62.9 \\
\bottomrule
\end{tabular*}
\end{table}

\subsection{Cumulative Incidence by Circuit}

Figures~\ref{fig:cif_circuit_dismissal}
and~\ref{fig:cif_circuit_settlement} present CIFs for the
four largest circuits. Circuit heterogeneity is substantial
and the patterns are not symmetric across outcomes---a
nuanced finding that a competing-risks framework reveals
and that standard single-outcome analyses would miss.

\begin{figure}[htb]
    \centering
    \begin{subfigure}[b]{0.49\textwidth}
        \centering
        \includegraphics[width=\textwidth]{./Figures/figure4_cif_circuit_dismissal.png}
        \caption{Dismissal Incidence by Circuit}
        \label{fig:cif_circuit_dismissal}
    \end{subfigure}
    \hfill
    \begin{subfigure}[b]{0.49\textwidth}
        \centering
        \includegraphics[width=\textwidth]{./Figures/figure5_cif_circuit_settlement.png}
        \caption{Settlement Incidence by Circuit}
        \label{fig:cif_circuit_settlement}
    \end{subfigure}
    \vskip 10pt
    \caption[Cumulative Incidence of Outcomes by Judicial Circuit]{Aalen-Johansen CIFs by circuit (top 4 by volume). Gray's test rejects equality for dismissal ($\text{stat} = 812.6, p < 0.001$) and settlement ($\text{stat} = 38.2, p < 0.001$).}
    \label{fig:cif_circuit}
\end{figure}

For dismissal (Figure~\ref{fig:cif_circuit_dismissal}),
the Second Circuit exhibits the highest long-run dismissal
incidence, reaching approximately 75\% by six years. The
Fifth Circuit leads in early dismissal---nearly 40\% by year
one---reflecting aggressive early motion practice in Texas
and Louisiana. The Third Circuit has the lowest overall
dismissal incidence, plateauing near 50\%.

The settlement picture is qualitatively different
(Figure~\ref{fig:cif_circuit_settlement}). The Third Circuit
has the highest settlement incidence, reaching approximately
30\% by year three. The Second Circuit has the lowest
settlement incidence of any major circuit, barely reaching
5\% by six years. This combination---low settlement
probability and high dismissal incidence---is the defining
empirical signature of the Southern District of New York.

\section{Baseline Cause-Specific Cox Models}

Table~\ref{tab:cox_baseline} presents baseline cause-specific
Cox models. Post-PSLRA cases have a 62\% lower settlement
hazard (HR $= 0.378$, 95\% CI: 0.337--0.424) and a 64\%
higher dismissal hazard (HR $= 1.638$, 95\% CI:
1.478--1.815). The Grambsch-Therneau test fails to reject
proportional hazards for settlement ($p = 0.12$) but
strongly rejects for dismissal ($p = 0.0001$), motivating
the piecewise analysis in Section~\ref{sec:piecewise}.

\begin{table}[htb]
\centering
\caption{Baseline Cause-Specific Cox Models (PSLRA Only)}
\label{tab:cox_baseline}
\renewcommand{\arraystretch}{1.2}
\begin{tabular*}{\textwidth}{@{\extracolsep{\fill}} l r c r r @{}}
\toprule
\textbf{Outcome} & \textbf{HR} & \textbf{95\% CI} & \textbf{$p$} & \textbf{PH $p$} \\
\midrule
Settlement & 0.378 & (0.337, 0.424) & $<$0.001 & 0.120 \\
Dismissal  & 1.638 & (1.478, 1.815) & $<$0.001 & 0.0001 \\
\bottomrule
\end{tabular*}
\end{table}

\subsection{Time-Varying PSLRA Effect on Dismissal Hazard}
\label{sec:piecewise}

Table~\ref{tab:piecewise} presents the piecewise constant
dismissal hazard model. The temporal decomposition reveals
a three-phase PSLRA effect. In the first year, post-PSLRA
cases face more than double the dismissal hazard of
pre-PSLRA cases (HR $= 2.18$): the PSLRA's heightened
pleading requirements and automatic discovery stay
concentrate dismissals at the pleading stage. In years one
to two, the PSLRA effect vanishes entirely (HR $= 0.922$,
$p = 0.444$): cases that survived the first-year gauntlet
face no differential dismissal risk. For cases surviving
beyond two years, a modest but significant PSLRA effect
re-emerges (HR $= 1.55$).

\begin{table}[htb]
\centering
\caption{PSLRA Effect on Dismissal Hazard by Time Period}
\label{tab:piecewise}
\renewcommand{\arraystretch}{1.2}
\begin{tabular*}{\textwidth}{@{\extracolsep{\fill}} l r c r @{}}
\toprule
\textbf{Period} & \textbf{HR} & \textbf{95\% CI} & \textbf{$p$} \\
\midrule
0--1 years & 2.180 & (1.850, 2.560) & $<$0.001 \\
1--2 years & 0.922 & (0.748, 1.140) & 0.444    \\
2$+$ years & 1.550 & (1.300, 1.840) & $<$0.001 \\
\bottomrule
\end{tabular*}
\end{table}

This three-phase pattern---strong early effect, null
middle-period effect, modest late re-emergence---is the
first formal empirical documentation of the PSLRA's
time-varying causal signature using a piecewise survival
framework on population-level data.

\section{Circuit Effects in Cause-Specific Cox Models}

Table~\ref{tab:cox_circuit} presents cause-specific Cox
models with PSLRA and circuit fixed effects. The Second
Circuit serves as the reference category.

\begin{table}[htb]
\centering
\caption{Cause-Specific Cox Models with Circuit Effects (Ref: Second Circuit)}
\label{tab:cox_circuit}
\renewcommand{\arraystretch}{1.2}
\begin{tabular*}{\textwidth}{@{\extracolsep{\fill}} l r c r c @{}}
\toprule
& \multicolumn{2}{c}{\textbf{Settlement}} & \multicolumn{2}{c}{\textbf{Dismissal}} \\
\cmidrule(lr){2-3} \cmidrule(lr){4-5}
\textbf{Covariate} & \textbf{HR} & \textbf{95\% CI} & \textbf{HR} & \textbf{95\% CI} \\
\midrule
Post-PSLRA    & 0.609 & (0.542, 0.684)
              & 1.846 & (1.664, 2.049) \\
First Circ.   & 8.52  & (6.26, 11.61)
              & 1.385 & (1.203, 1.595) \\
Third Circ.   & 16.57 & (13.53, 20.30)
              & 1.178 & (1.054, 1.317) \\
Fourth Circ.  & 1.765 & (1.397, 2.230)
              & 0.262 & (0.228, 0.301) \\
Fifth Circ.   & 7.519 & (5.929, 9.536)
              & 1.483 & (1.343, 1.639) \\
Sixth Circ.   & 15.89 & (12.62, 20.02)
              & 0.726 & (0.607, 0.867) \\
Seventh Circ. & 9.587 & (7.570, 12.14)
              & 1.251 & (1.109, 1.412) \\
Eighth Circ.  & 13.60 & (10.70, 17.30)
              & 0.913 & (0.788, 1.059) \\
Ninth Circ.   & 7.359 & (6.188, 8.752)
              & 1.103 & (1.039, 1.171) \\
Tenth Circ.   & 8.290 & (6.529, 10.52)
              & 0.805 & (0.703, 0.923) \\
Eleventh Circ.& 8.265 & (6.438, 10.61)
              & 1.607 & (1.447, 1.784) \\
\bottomrule
\end{tabular*}
\end{table}

The most striking finding is the extreme magnitude of circuit
effects on settlement hazard. Every other circuit has a
settlement hazard 1.8 to 16.6 times higher than the Second
Circuit's. The Second Circuit is, by a wide margin, the
jurisdiction in which securities class actions are least
likely to settle at any given duration. On the dismissal
side, the Fourth Circuit stands out in the opposite
direction: its dismissal hazard is only 26\% of the Second
Circuit's (HR $= 0.262$), the lowest of any circuit.

\section{Extended Covariate Models}

The extended analysis uses 12,866 cases with complete data
on origin, MDL flag, and jurisdiction. Statutory basis
coverage is 97.5\%; a missingness indicator preserves the
full estimation sample.

\subsection{Extended Cause-Specific Cox Models}

Tables~\ref{tab:cox_ext_settle}
and~\ref{tab:cox_ext_dismiss} present the extended
cause-specific Cox models. The PSLRA effects are robust to
the inclusion of all extended covariates: the settlement HR
remains 0.617 and the dismissal HR increases slightly to
1.941. The PSLRA structural break is not an artifact of
compositional differences in origin, MDL status, or
statutory basis across regimes.

\begin{table}[htbp]
\centering
\caption{Extended Cause-Specific Cox Model: Settlement ($n = 12{,}866$)}
\label{tab:cox_ext_settle}
\renewcommand{\arraystretch}{1.1} % Slightly tighter for long lists, but with better padding
\begin{tabular*}{\textwidth}{@{\extracolsep{\fill}} l r c @{}}
\toprule
\textbf{Covariate} & \textbf{HR} & \textbf{95\% CI} \\
\midrule
Post-PSLRA            & 0.617 & (0.548, 0.695) \\
\addlinespace[0.5em]
\textit{Circuit (Ref: Second Circuit)} & & \\
Third Circuit         & 11.62 & (9.414, 14.340) \\
Sixth Circuit         & 11.49 & (9.027, 14.612) \\
Eighth Circuit        & 10.08 & (7.880, 12.889) \\
Ninth Circuit         &  5.20 & (4.338,  6.222) \\
Fifth Circuit         &  5.84 & (4.581,  7.442) \\
First Circuit         &  6.04 & (4.413,  8.259) \\
Seventh Circuit       &  6.88 & (5.388,  8.772) \\
Tenth Circuit         &  5.99 & (4.630,  7.737) \\
Eleventh Circuit      &  5.56 & (4.300,  7.194) \\
Fourth Circuit        &  2.93 & (2.058,  4.172) \\
\addlinespace[0.5em]
\textit{Case Characteristics} & & \\
Origin: Removed       &  1.44 & (0.926,  2.226) \\
Origin: MDL           &  0.351 & (0.252, 0.489) \\
Origin: Other         &  1.57 & (1.336,  1.834) \\
MDL Flag              &  0.262 & (0.181, 0.380) \\
Federal Question      &  4.37 & (1.635, 11.672) \\
\addlinespace[0.5em]
\textit{Statutory Basis (Ref: 10(b))} & & \\
Section 11            &  1.02 & (0.861,  1.202) \\
Other/Both Statutory  &  1.01 & (0.882,  1.146) \\
\bottomrule
\end{tabular*}
\end{table}

\begin{table}[htbp]
\centering
\caption{Extended Cause-Specific Cox Model: Dismissal ($n = 12{,}866$)}
\label{tab:cox_ext_dismiss}
\renewcommand{\arraystretch}{1.1}
\begin{tabular*}{\textwidth}{@{\extracolsep{\fill}} l r c @{}}
\toprule
\textbf{Covariate} & \textbf{HR} & \textbf{95\% CI} \\
\midrule
Post-PSLRA            & 1.941 & (1.748, 2.157) \\
\addlinespace[0.5em]
\textit{Circuit (Ref: Second Circuit)} & & \\
Third Circuit         & 0.984 & (0.879, 1.101) \\
Sixth Circuit         & 0.632 & (0.526, 0.760) \\
Eighth Circuit        & 0.808 & (0.696, 0.938) \\
Ninth Circuit         & 0.907 & (0.855, 0.964) \\
Fifth Circuit         & 1.298 & (1.174, 1.434) \\
First Circuit         & 1.158 & (1.006, 1.334) \\
Seventh Circuit       & 0.909 & (0.800, 1.034) \\
Tenth Circuit         & 0.714 & (0.618, 0.824) \\
Eleventh Circuit      & 1.265 & (1.139, 1.405) \\
Fourth Circuit        & 0.597 & (0.497, 0.717) \\
\addlinespace[0.5em]
\textit{Case Characteristics} & & \\
Origin: Removed       & 0.783 & (0.615, 0.998) \\
Origin: MDL           & 0.227 & (0.186, 0.276) \\
Origin: Other         & 0.975 & (0.889, 1.069) \\
MDL Flag              & 0.274 & (0.242, 0.310) \\
Federal Question      & 2.290 & (1.545, 3.395) \\
\addlinespace[0.5em]
\textit{Statutory Basis (Ref: 10(b))} & & \\
Section 11            & 1.121 & (1.041, 1.207) \\
Other/Both Statutory  & 1.015 & (0.947, 1.088) \\
\bottomrule
\end{tabular*}
\end{table}

MDL consolidation has a large and symmetric effect on both
outcomes. Cases assigned to an MDL docket have a 74\% lower
settlement hazard (HR $= 0.262$) and a 73\% lower dismissal
hazard (HR $= 0.274$). This bilateral suppression indicates
that MDL consolidation dramatically prolongs case pendency by
suspending normal resolution dynamics while coordinating
discovery and bellwether proceedings across related cases.
Federal question jurisdiction is associated with higher
hazards for both outcomes. Section~11 cases have a modestly
higher dismissal hazard than Section~10(b) cases (HR $=
1.121$, 95\% CI: 1.041--1.207), with no significant
difference in settlement hazard.

Grambsch-Therneau global tests reject proportional hazards
in both extended models (settlement: $\chi^2 = 102.9$,
$p < 0.001$; dismissal: $\chi^2 = 436.9$, $p < 0.001$),
driven primarily by circuit, origin, MDL flag, and statutory
basis. The extended Cox estimates should be interpreted as
average effects over the study period. Full time-varying
coefficient models are deferred to the final thesis.

\subsection{Extended Fine-Gray Subdistribution Models}

Table~\ref{tab:fg_comparison} compares cause-specific and
Fine-Gray subdistribution hazard ratios for selected
covariates, illustrating the substantively important
divergence between the two frameworks.

\begin{table}[htbp]
\centering
\caption{Cause-Specific vs.\ Subdistribution Hazard Ratios: Selected Covariates, Extended Models}
\label{tab:fg_comparison}
\renewcommand{\arraystretch}{1.3} % More vertical space for clarity
\begin{tabular*}{\textwidth}{@{\extracolsep{\fill}} l l r r @{}}
\toprule
\textbf{Outcome} & \textbf{Covariate} & \textbf{CS HR} & \textbf{FG SHR} \\
\midrule
Settlement & Post-PSLRA & 0.617 & 0.411 \\
           & MDL Flag   & 0.262 & 0.963 \\
\midrule
Dismissal  & Post-PSLRA & 1.941 & 2.050 \\
           & MDL Flag   & 0.274 & 0.500 \\
\bottomrule
\end{tabular*}
\begin{minipage}{\textwidth}
\vskip 4pt
\footnotesize \textit{Note:} CS = cause-specific hazard ratio; FG = Fine-Gray subdistribution hazard ratio. CS models treat competing events as censored, while FG models account for the reduced "at-risk" pool.
\end{minipage}
\end{table}

For PSLRA, Fine-Gray subdistribution hazard ratios are more
extreme than cause-specific ratios for both outcomes. The
amplification occurs because the PSLRA increases dismissals,
reducing the pool of cases available to settle; the
subdistribution model captures this indirect pathway while
the cause-specific model holds the competing risk fixed.

The MDL comparison is more analytically revealing. Despite
the dramatically lower cause-specific settlement hazard (HR
$= 0.262$), the Fine-Gray subdistribution effect on
settlement is essentially null (SHR $= 0.963$). The
reconciliation is that MDL cases also have very low dismissal
hazard: because MDL cases rarely get dismissed, the pool
available to settle remains large throughout, and cumulative
settlement probability is not substantially lower than for
non-MDL cases despite the lower instantaneous rate. This
result is invisible to single-outcome analyses.

\section{PSLRA $\times$ Circuit Interaction}

Likelihood ratio tests comparing the additive
circuit-PSLRA specification against the interaction model
reject the null of a uniform PSLRA effect: settlement
($\chi^2 = 13.46$, df $= 5$, $p = 0.019$) and dismissal
($\chi^2 = 25.05$, df $= 5$, $p = 0.0001$). The PSLRA did
not operate uniformly across jurisdictions.

For settlement, all interaction terms are below 1. Because the 
baseline PSLRA effect (referenced to the Second Circuit) is 
already below 1, interaction multipliers below 1 indicate that 
the PSLRA's suppression of settlement hazard was \emph{amplified} 
in other circuits. In other words, the Second Circuit experienced 
the weakest absolute PSLRA-induced reduction in settlement hazard. 
This is consistent with the Second Circuit already possessing a 
highly restrictive baseline settlement environment, leaving less 
room for the PSLRA to further suppress settlement probabilities 
relative to more plaintiff-friendly venues. 

For dismissal, the most striking interaction is in the Eleventh 
Circuit (interaction HR $= 0.426$), indicating strong attenuation 
of the PSLRA's baseline positive effect on dismissals. The Ninth 
Circuit shows a slight amplification (HR $= 1.052$). These patterns 
are consistent with documented variation in judicial application of 
the ``strong inference'' scienter pleading standard across circuits 
\citep{johnson2001, choi2007}.

\section{Robustness Checks}

Figure~\ref{fig:robustness} and Table~\ref{tab:robustness}
present PSLRA hazard ratios from six distinct Cox models.
Across every specification, settlement HRs are below 1 and
dismissal HRs are above 1. The direction of the PSLRA effect
is not an artifact of any particular coding or sampling
choice.

\begin{figure}[htb]
    \centering
    \includegraphics[width=\linewidth]{./Figures/figure6_robustness_hr.png}
    \vskip 5pt
    \caption[PSLRA Hazard Ratios Across Robustness Specifications]{PSLRA hazard ratios across robustness specifications. Each row is a separate baseline Cox model. Green triangles = settlement HR; red circles = dismissal HR. Dashed line at HR $= 1$ (null). Note the log scale on the horizontal axis.}
    \label{fig:robustness}
\end{figure}

\begin{table}[htbp]
\centering
\caption{PSLRA Hazard Ratios Across Robustness Specifications}
\label{tab:robustness}
\renewcommand{\arraystretch}{1.2}
\begin{tabular*}{\textwidth}{@{\extracolsep{\fill}} l l r c r @{}}
\toprule
\textbf{Specification} & \textbf{Outcome} & \textbf{HR} & \textbf{95\% CI} & \textbf{$p$} \\
\midrule
Scheme A: Primary       & Settlement & 0.378 & (0.337, 0.424) & $<$0.001 \\
                        & Dismissal  & 1.638 & (1.478, 1.815) & $<$0.001 \\
\addlinespace[0.5em]
Scheme B: Code 12 = S   & Settlement & 0.696 & (0.630, 0.769) & $<$0.001 \\
                        & Dismissal  & 1.470 & (1.310, 1.650) & $<$0.001 \\
\addlinespace[0.5em]
Scheme C: Codes 12+5 = S & Settlement & 0.691 & (0.627, 0.761) & $<$0.001 \\
                        & Dismissal  & 1.470 & (1.310, 1.650) & $<$0.001 \\
\addlinespace[0.5em]
Excl. post-2020         & Settlement & 0.377 & (0.336, 0.423) & $<$0.001 \\
                        & Dismissal  & 1.570 & (1.420, 1.740) & $<$0.001 \\
\addlinespace[0.5em]
Second Circuit only     & Settlement & 0.794 & (0.396, 1.590) & 0.515    \\
                        & Dismissal  & 1.680 & (1.300, 2.180) & $<$0.001 \\
\addlinespace[0.5em]
Ninth Circuit only      & Settlement & 0.771 & (0.629, 0.944) & 0.012    \\
                        & Dismissal  & 2.340 & (1.940, 2.830) & $<$0.001 \\
\bottomrule
\end{tabular*}
\end{table}

Three observations warrant emphasis. First, Scheme~B and~C
estimates for settlement (HR $\approx 0.69$--$0.70$) are
larger than Scheme~A (HR $= 0.378$), as expected when
Code~12 cases are reclassified as settlements. Even so, the
effect remains large and highly significant under all liberal
coding assumptions. Second, excluding post-2020 cases has
virtually no effect on estimates, indicating that recently
filed cases with high censoring rates do not distort the
main results. Third, the Second Circuit sub-model yields a
non-significant PSLRA effect on settlement hazard ($p =
0.515$): within the Second Circuit, the PSLRA did not
significantly alter settlement dynamics, though it did
significantly increase the dismissal hazard (HR $= 1.68$).
The Ninth Circuit, by contrast, shows a significant PSLRA
effect on both outcomes, with a dismissal HR of 2.34 that
substantially exceeds the pooled estimate.

\section{Random Survival Forest: Variable Importance
and Predictive Performance}

\subsection{Variable Importance}

Figure~\ref{fig:rsf_vimp} presents permutation-based VIMP
scores from the final RSF (200 trees, $m_{\text{try}} = 2$,
selected by 5-fold cross-validation with optimal OOB error
of 0.408).

\begin{figure}[htb]
    \centering
    \includegraphics[width=\linewidth]{./Figures/figure7_rsf_vimp.png}
    \vskip 5pt
    \caption[RSF Variable Importance by Outcome]{Random Survival Forest variable importance (VIMP) by outcome. VIMP represents the increase in OOB prediction error when the variable is randomly permuted. Left: dismissal; right: settlement. RSF: 200 trees, $m_{\text{try}} = 2$.}
    \label{fig:rsf_vimp}
\end{figure}

Circuit identity is the single most important predictor for
both outcomes (VIMP: 0.087 for settlement, 0.154 for
dismissal). This nonparametric confirmation of circuit
dominance reinforces the parametric findings in
Sections~5.3 and~5.4. The PSLRA indicator ranks second for
dismissal (VIMP $= 0.130$), confirming the reform's
importance as a predictor of dismissal outcomes. For
settlement, MDL flag ranks second (VIMP $= 0.074$),
reflecting the bilateral hazard suppression documented in
Section~5.5. Statutory basis has the lowest VIMP, consistent
with its near-null coefficients in the extended Cox models.

\subsection{Model Comparison: Discrimination Performance}

Table~\ref{tab:performance} and Figure~\ref{fig:auc_comparison}
compare predictive discrimination across models on the
held-out 30\% test set ($n = 3{,}861$).
\begin{table}[htbp]
\centering
\caption{Model Performance: Held-Out Test Set ($n = 3{,}861$)}
\label{tab:performance}
\renewcommand{\arraystretch}{1.2}
\begin{tabular*}{\textwidth}{@{\extracolsep{\fill}} l r r r r @{}}
\toprule
\textbf{Model} & \textbf{C-Idx} & \textbf{AUC@1} & \textbf{AUC@3} & \textbf{AUC@5} \\
\midrule
Cox (Settlement)       & 0.735 & 0.746 & 0.741 & 0.836 \\
Cox (Dismissal)        & 0.602 & 0.606 & 0.639 & 0.772 \\
\addlinespace[0.5em]
Fine-Gray (Settlement) & 0.735 & ---   & ---   & ---   \\
Fine-Gray (Dismissal)  & 0.602 & ---   & ---   & ---   \\
\addlinespace[0.5em]
RSF (Settlement)       & 0.549 & 0.538 & 0.598 & 0.768 \\
RSF (Dismissal)        & 0.540 & 0.554 & 0.622 & 0.735 \\
\bottomrule
\end{tabular*}
\begin{minipage}{\textwidth}
\vskip 4pt
\footnotesize \textit{Note:} Fine-Gray and Cox share the same linear predictor under the \texttt{finegray} augmented-data implementation; C-indices are therefore identical. Separate Fine-Gray AUC computation and Integrated Brier Score (IBS) are deferred to the final thesis.
\end{minipage}
\end{table}

\begin{figure}[htb]
    \centering
    \includegraphics[width=\linewidth]{./Figures/figure8_auc_comparison.png}
    \vskip 5pt
    \caption[Time-Dependent AUC for Cox and RSF Models]{Time-dependent AUC for Cox and RSF at 1, 2, 3, and 5-year prediction horizons. Solid lines = Cox; dashed = RSF. Green = settlement; red = dismissal. Dotted line at AUC $= 0.5$ indicates no-skill performance.}
    \label{fig:auc_comparison}
\end{figure}

The Cox model substantially outperforms the current proof-of-concept 
RSF implementation across every metric and outcome. For settlement, 
the Cox C-index is 0.735 versus RSF's 0.549, a gap of 18.6 percentage 
points. For dismissal, Cox C-index is 0.602 versus RSF's 0.540. The 
AUC comparison confirms this: Cox settlement AUC is consistently
0.74--0.84 across horizons, while RSF settlement AUC ranges
from 0.54 to 0.77. 

The near-random discrimination of the RSF (C-index $\approx 0.54$) 
is highly characteristic of an underpowered ensemble. Because this 
interim implementation utilizes only 200 trees and a limited 
$m_{\text{try}}$ grid search over highly censored data, the trees 
likely lack the depth and variance necessary to stabilize the 
cumulative hazard estimates. While the Cox model's linear predictor 
currently captures the dominant main effects (circuit and PSLRA regime) 
effectively, a definitive conclusion regarding the value of nonlinear 
ML methods in this domain must await the fully scaled, computationally 
tuned RSF ensemble planned for the final thesis.

\section{Summary of Results}

Our analysis yields six main findings. First, the PSLRA
fundamentally restructured securities litigation outcomes,
with effects robust to all disposition coding schemes,
temporal restrictions, and extended covariate controls.
Second, the PSLRA's impact on dismissal hazard is
concentrated in the first year (HR $= 2.18$), vanishes in
year two (HR $= 0.922$, $p = 0.444$), and re-emerges
modestly thereafter (HR $= 1.55$). Third, circuit geography
is the dominant predictor of case outcomes in both
parametric and nonparametric frameworks, with the Second
Circuit exhibiting anomalously low settlement hazard and
the highest dismissal incidence. Fourth, the PSLRA's effects
are heterogeneous across circuits, with the reform's settlement 
suppression being weakest in the Second Circuit and its dismissal 
amplification strongly attenuated in the Eleventh Circuit. 
Fifth, MDL consolidation suppresses both cause-specific hazards, 
but Fine-Gray analysis reveals no lower cumulative settlement 
probability---a result invisible to single-outcome analyses. 
Sixth, the semi-parametric Cox model currently achieves strong 
out-of-sample discrimination (Settlement C-Index = 0.735), 
establishing a rigorous baseline that the final scaled machine 
learning models must overcome.